% Market Structure and Information Acquisition

We next examine whether the informativeness of the lending signal varies with the structure of the lending market. If lending conveys private information about short-selling demand, the signal should be more informative in markets where lending is concentrated---i.e., where each individual lending transaction is more likely to reflect informed borrowing. We test this hypothesis by interacting our baseline lending and disclosure variables with measures of lending market concentration.

Table~\ref{table_market_structure} presents the results, all estimated with the full set of saturated fixed effects (security$\times$time, fund$\times$time, fund$\times$security). Column~(1) reproduces the baseline for reference: $OnLoan_{i,j,t-1} = -0.4226$ ($t = -2.35$) with a disclosure interaction of $-0.8633$ ($t = -2.32$).

We consider five measures of market concentration. First, at the stock level, we use the Herfindahl-Hirschman Index of lending concentration ($HHI_{i,t-1}$) and the inverse of the number of active lending agents ($1/ActiveAgents_{i,t-1}$). Second, at the fund level, we use the fund family's lending concentration ($HHI_{f,t-1}$), the inverse of the number of stocks in the fund's portfolio ($1/Stocks_{f,t-1}$), and the inverse of the number of stocks on loan ($1/Stocks\_OnLoan_{f,t-1}$).

The triple-interaction results support the information concentration hypothesis. The coefficient on $OnLoan \times 1/ActiveAgents \times Disclosure$ is $-9.2142$ ($t = -2.41$): when fewer agents participate in the lending market for a stock, the lending signal is substantially more informative. Similarly, fund-family concentration amplifies the effect: $OnLoan \times HHI_f \times Disclosure = -3.0048$ ($t = -1.94$). The strongest result comes from the inverse of stocks on loan: $OnLoan \times 1/Stocks\_OnLoan \times Disclosure = -8.1425$ ($t = -3.96$). When a fund lends fewer distinct securities, each lending position carries more informational weight, and the trading response is correspondingly larger.

These heterogeneity results reinforce the information acquisition interpretation. The trading response to lending is not uniform across all funds and stocks; rather, it is amplified precisely in the settings where the information content of individual lending transactions is highest. This pattern is difficult to reconcile with mechanical explanations and is instead consistent with a model in which fund managers extract and act upon informative signals embedded in the borrowing demand they observe through their lending operations.
